\section{Introduction}
\label{sec:introduction}

% Motivation: why solving polynomial systems is important, and what we want to find
Systems of multivariate polynomial equations appear whenever a symbolic model
contains nonlinear algebraic constraints.
Applications include robotics, coding theory, optimization, statistics, and physics \cite{sturmfels2002solving}.
In many of these problems, an isolated solution is not enough.
One needs the dimension of the solution set, its irreducible components, or exact relations that hold on every solution.
These are questions about the variety and its defining ideal, rather than about a single point \cite{cox2015ideals,kreuzer2000computational}.

% (minor) give an example here, of where a particular solution is less interesting than the general family

% History of the field and various methods
Elimination by resultants provided an early systematic way to attack polynomial
systems, but the introduction of Gr\"obner bases gave a common algorithmic
framework for elimination, ideal membership, dimension, and related tasks.
Buchberger developed this framework in his 1965 thesis \cite{buchberger2006thesis}.
Later algorithms changed how the expensive reductions are organized.
F4 replaces individual polynomial reductions with sparse linear algebra, while F5 uses signatures to avoid many reductions to zero \cite{faugere1999f4,faugere2002f5}.
Primary-decomposition algorithms use Gr\"obner bases and factorization to separate components \cite{gianni1988primary}.
Numerical algebraic geometry takes a different route: homotopy continuation and witness sets provide isolated roots numerically \cite{sommese2001numerical,sommese2005numerical}.

% Mention recent works in a sentence; emphasize that the field is still active
Recent learning-based work has concentrated on these choices.
Peifer et al.\ trained a reinforcement-learning policy for pair selection in
Buchberger's algorithm \cite{peifer2020learning}. Kera et al.\ trained
transformers to compute Gr\"obner bases for a restricted class of
zero-dimensional radical ideals \cite{kera2024learning}. Both results concern
controlled input distributions; a large, general ideal still requires an
exact symbolic computation.

% What's the gap?
% A statement on complexity and limit of what is feasible.
% Despite a lot of progress... many problems with polynomial equations remain inaccessible, including ones of particular relevance in quantum field theory and more.
% There are also many situations of interest where ceratin domain-... assumptions can be used to ...

% Here we (expand on the abstract)

% Technical description of the method
Our method starts with that factor structure. Many equations in our target
systems factor. If
$p_i=\prod_j f_{ij}$, then every solution of $p_i=0$ makes at least one factor
$f_{ij}$ vanish. Selecting factors that meet every equation gives a hitting set
of the factor family. It is enough to enumerate the minimal hitting sets, and
each one defines a smaller polynomial system. We solve these systems separately.
The union of their zero sets is exactly the zero set of the original input, so
the split does not discard solutions or rely on a genericity assumption.

% see what of the below is needed above - most of it goes to later sections
The framework uses UNSAT core learning to prune hitting sets and the accelerate
computation.
The method does not require a particular branch solver or core-extraction procedure.

Some factor branches can still be too expensive for the chosen subsolver.
After a timeout, we branch again splitting a selected variable into the cases $x=0$ and $x\neq 0$.
A domain expert may have a good reason to choose a different splitting which will yield faster computation.

% keep things we're proud of, like nontrivial proofs
We prove that our methods preserve the complete solution set for any branch solver satisfying the stated conditions.
We compare the resulting implementation with direct Gr\"obner-basis computation, symbolic decomposition, and numerical homotopy on generated benchmarks.
We then apply it to two systems from quantum field theory for which the standard methods in our comparison do not complete. The output is an exact collection of branch descriptions for the full solution set, including
positive-dimensional families.
